\section{Objective-specific UCB and Thompson-style policies}
Let $b^*=\argmax_i\mu_i$. We distinguish three cumulative comparators:
\begin{align}
R_{\rm oracle}(T)&=\sum_{t\le T}(\max_iX_{i,t}-X_{a_t,t}),\nonumber\\
R_\mu(T)&=\sum_{t\le T}(\mu_{b^*}-\mu_{a_t}),\label{eq:three-regrets}\\
R_{\rm fixed}(T)&=\sum_{t\le T}(X_{b^*,t}-X_{a_t,t}).\nonumber
\end{align}
The first two are nonnegative pathwise. The third may be negative: adaptive actions can exploit positive predictable deviations. Since the deviation of the fixed arm has mean zero,
\begin{equation}
\E R_{\rm fixed}(T)=\E R_\mu(T)-\sum_{t\le T}\E z_{a_t,t}.
\end{equation}
Even the full-current-state oracle can have positive mean pseudo-regret. Terminal PCS instead compares the recommendation's permanent mean with $\mu_{b^*}$.

\subsection{A shared inference engine}
For observations strictly before decision $t$, write their design, covariance and readings as $M,\Xi,Y$. Define
\begin{align}
V&=(H+M^\top\Xi^{-1}M)^{-1},&m&=VM^\top\Xi^{-1}Y,\nonumber\\
K&=\Cov(z_t,\hbox{selected past deviations}),&A&=K\Xi^{-1},\nonumber\\
W&=I-AM,&f&=AY+Wm,\label{eq:forecast-working}\\
P&=G_0-K\Xi^{-1}K^\top,&\Sigma&=P+WVW^\top.\nonumber
\end{align}
At the true fixed mean, the conditional expected field is $AY+W\mu$, with conditional noise covariance $P$. The matrix $\Sigma$ is a working Gaussian uncertainty for randomized decisions; $V$ remains inverse penalized information, not the repeated-noise sampling covariance. A Gaussian integration of the penalized likelihood produces these same working quantities, without specifying a random physical mean population. Correct frequentist uncertainty uses~\eqref{eq:confidence}, including its regularization allowance.

The selected-history implementation updates $\Xi^{-1}$ by its Schur complement. It is algebraically equivalent to the affine state filter, retains the full AR likelihood, and uses $O(T^2)$ history storage, $O(Nt^2)$ all-arm forecast work, and $O(t^2+N^2)$ selected-update work. The lag-state alternative has dimension $\sum_i p_i$. Neither is a constant-cost solution as all dimensions grow.

\subsection{Permanent-mean UCB and a regret bound}
Mean UCB chooses
\begin{equation}
a_t=\argmax_i\{m_i+c_t\sqrt{V_{ii}}\}.
\label{eq:mean-ucb}
\end{equation}
The certified version uses $c_t=\beta_{t-1}(\delta)$. The practical version uses $c_t=.7\sqrt{2\log(t+1)}$ for one-based decisions. The practical scale has no automatic coverage claim.

After initially measuring each arm, force cyclic probes at square calendar ages: if $t-N=k^2$ for $k\ge1$, select $(k-1)\bmod N$. Other rounds use the policy. Thus the forced count satisfies $F_T\le N+\sqrt T$, and each arm receives $\Omega(\sqrt T/N)$ forced samples eventually.

\begin{theorem}[Conservative stationary mean-regret guarantee]
Under the assumptions of the confidence theorem and the spectral covariance bounds~\eqref{eq:information-bounds}, certified mean UCB with the preceding probes obeys, with probability at least $1-\delta$,
\begin{equation}
\boxed{R_\mu(T)\le 2BF_T+4\beta_T\sqrt{v_+NT}.}
\label{eq:mean-regret-bound}
\end{equation}
Moreover,
\[
\beta_T\le\sqrt{N\log(1+T/(\alpha v_-))+2\log(1/\delta)}
+\sqrt{\alpha B^2+\lambda S^2}.
\]
This bound is sublinear for fixed $N$ and fixed stable dynamics. It does not claim sharp graph or temporal constants.
\end{theorem}
\begin{proof}
On the simultaneous confidence event, a non-forced UCB action has gap at most $2\beta_{t-1}\sqrt{V_{a_ta_t}}$. After initial coverage,~\eqref{eq:information-bounds} implies $V_{ii}\le v_+/N_i(t-1)$. Summing inverse square-root counts over ordinary rounds is bounded by $2\sum_i\sqrt{N_i(T)}\le2\sqrt{NT}$. Forced actions cost at most $2B$ each. The determinant bound follows from $H\succeq\alpha I$ and the upper information bound.
\end{proof}
An expectation bound adds $2BT\delta$; choosing $\delta=1/T$ in a horizon-specific policy makes that remainder bounded. The experiment uses $\delta=.05$ for the certified policy. Because minimum forced counts grow as $\sqrt T$, the confidence widths vanish as $O(\sqrt{\log T}/T^{1/4})$ for fixed dimensions. Applying the confidence theorem with arbitrarily small failure probabilities gives consistent mean estimates and PCS tending to one at a positive fixed gap for every correctly specified joint policy with these probes. This asymptotic statement is not a finite-budget PCS guarantee.
